\documentclass[12pt]{article}

% Codificación
\usepackage[utf8]{inputenc}
\usepackage[T1]{fontenc}
\usepackage[spanish]{babel}

% Matemáticas
\usepackage{amsmath}
\usepackage{amssymb}
\usepackage{amsfonts}

% Algoritmos
\usepackage{algorithm}
\usepackage{algpseudocode}

% Figuras
\usepackage{graphicx}

% Tablas
\usepackage{booktabs}

% Control de flotantes
\usepackage{float}

% Márgenes
\usepackage{geometry}

% Links
\usepackage{hyperref}

% Diagramas
\usepackage{tikz}
\usetikzlibrary{arrows.meta, positioning, fit}

\title{Evaluacion04}
\date{March 2026}

\begin{document}

% --- Portada ---
\begin{titlepage}
\begin{center}

\vspace*{1.5cm}

{\Large \textbf{UNIVERSIDAD POLITÉCNICA DE CHIAPAS}}\\[0.4cm]
{\large MAESTRÍA EN CIENCIAS E INNOVACIÓN TECNOLÓGICA}

\vspace{2cm}

{\large \textbf{ASIGNATURA}}\\[0.2cm]
{\large LÓGICA Y PROGRAMACIÓN}

\vspace{1.8cm}

{\Large \textbf{EVALUACIÓN 04}}

\vspace{1.8cm}

{\large \textbf{TRABAJO GRUPAL}}\\[0.2cm]

\vspace{2cm}

{\large \textbf{DOCENTE}}\\[0.2cm]
{\large DR. HORACIO IRÁN SOLÍS CISNEROS}

\vfill

{\large Tuxtla Gutiérrez, Chiapas\\
05 de marzo de 2026}

\end{center}
\end{titlepage}

\section{Introducción}

\subsection{Contexto}

El crecimiento del tráfico vehicular en entornos urbanos representa uno de los principales desafíos para la movilidad en ciudades modernas. En el marco del concepto de \textit{ciudades inteligentes}, los sistemas de control de tráfico buscan mejorar la eficiencia del flujo vehicular mediante el uso de tecnologías de información, sensores y algoritmos de optimización que permitan ajustar dinámicamente los tiempos de los semáforos.

En muchas implementaciones tradicionales, la sincronización de semáforos se realiza mediante algoritmos secuenciales que procesan cada intersección de manera independiente y en orden. Este enfoque simplifica la implementación del sistema, pero introduce limitaciones cuando el número de intersecciones aumenta o cuando se requiere procesar información en tiempo cercano al tiempo real.

En este trabajo se analiza inicialmente el comportamiento de un sistema de sincronización de semáforos bajo un modelo secuencial, denominado \textbf{Estado \(P\)}, con el objetivo de identificar formalmente sus limitaciones computacionales y los cuellos de botella que afectan su rendimiento.

A partir de este análisis, el trabajo plantea posteriormente la construcción de un modelo alternativo, denominado \textbf{Estado \(Q\)}, en el cual el procesamiento de intersecciones puede realizarse mediante mecanismos de concurrencia o paralelización. La comparación entre ambos estados permitirá evaluar el impacto de la concurrencia en el rendimiento y escalabilidad del sistema de sincronización de semáforos.

\subsection{Problema}

Considérese una red urbana compuesta por un conjunto de intersecciones controladas por semáforos:

\[
I = \{i_1, i_2, ..., i_n\}
\]

Cada intersección requiere ejecutar una serie de operaciones para determinar el estado del semáforo en cada ciclo de control.

\begin{itemize}
\item adquisición de datos de tráfico
\item cálculo de densidad vehicular
\item determinación del tiempo óptimo del semáforo
\item actualización del estado del semáforo
\end{itemize}

Formalmente, para cada intersección \(i\) se define la función de procesamiento:

\[
F(i) = read(i) \rightarrow compute(i) \rightarrow update(i)
\]

\subsection{Estado P: Sistema Secuencial}

En el estado inicial del sistema, denominado \(P\), el procesamiento de todas las intersecciones se realiza de forma estrictamente secuencial.

\[
P = \sum_{i=1}^{n} F(i)
\]

donde:

\[
F(i) = read(i) + compute(i) + update(i)
\]

El tiempo total de ejecución puede aproximarse como:

\[
T_P = \sum_{i=1}^{n} (t_r + t_c + t_u)
\]

donde

\begin{itemize}
\item \(t_r\) tiempo de lectura
\item \(t_c\) tiempo de cómputo
\item \(t_u\) tiempo de actualización
\end{itemize}

Por lo tanto,

\[
T_P = n(t_r + t_c + t_u)
\]

lo que implica una complejidad temporal:

\[
T_P = O(n)
\]

\subsection{Hipótesis de Complejidad}

El análisis de complejidad permite estimar el comportamiento del sistema conforme aumenta el número de intersecciones en la red de tráfico.

\subsubsection{Complejidad por intersección}

Cada intersección ejecuta el pipeline:

\[
F(i) = read(i); compute(i); update(i)
\]

Cada operación tiene un costo constante, por lo que:

\[
F(i) = O(1)
\]

\subsubsection{Complejidad del sistema completo}

El sistema procesa \(n\) intersecciones de forma secuencial:

\[
P = F(i_1); F(i_2); \dots ; F(i_n)
\]

Por lo tanto:

\[
T_P = \sum_{i=1}^{n} F(i)
\]

lo cual produce:

\[
T_P = O(n)
\]

\subsubsection{Hipótesis}

Se plantea la siguiente hipótesis de complejidad:

\textbf{Hipótesis:}

Si el sistema de sincronización de semáforos se ejecuta de forma secuencial, entonces el tiempo total de ejecución crecerá linealmente con el número de intersecciones procesadas.

Esto implica que, en redes de tráfico grandes, el sistema puede presentar limitaciones de escalabilidad y tiempos de respuesta elevados.

\subsection{Cuello de botella}

El principal cuello de botella del sistema secuencial radica en que cada intersección debe esperar a que la anterior complete su procesamiento.

Esto implica:

\begin{itemize}
\item ausencia de paralelismo
\item bloqueo del flujo de procesamiento
\item crecimiento lineal del tiempo de ejecución
\end{itemize}

\subsection{Pregunta de investigación}

A partir de este análisis surge la siguiente pregunta:

\textbf{¿Cómo puede optimizarse el procesamiento del sistema para reducir los tiempos de ejecución del modelo secuencial sin comprometer su corrección funcional?}

\subsection{Especificación formal mediante tripletas de Hoare}

\[
\{traffic\_data(i)\} \ read(i) \ \{raw\_data(i)\}
\]

\[
\{raw\_data(i)\} \ compute(i) \ \{timing(i)\}
\]

\[
\{timing(i)\} \ update(i) \ \{light\_state(i)\}
\]

El pipeline completo se expresa como:

\[
\{traffic\_data(i)\} \ F(i) \ \{light\_state(i)\}
\]

y el sistema completo:

\[
\{TrafficNetwork\} \ P \ \{UpdatedTrafficLights\}
\]

\subsection{Diagrama del pipeline (composición secuencial)}

\begin{figure}[H]
\centering
\begin{tikzpicture}[
  font=\small,
  node distance=6mm and 10mm,
  block/.style={rectangle, draw, rounded corners, align=center, minimum height=9mm, minimum width=22mm},
  inner/.style={rectangle, draw, rounded corners, align=center, minimum height=6mm, minimum width=16mm},
  arrow/.style={-Latex, thick},
  bott/.style={draw, dashed, rounded corners, inner sep=4mm}
]

\node[inner] (r1) {$read(i_1)$};
\node[inner, below=of r1] (c1) {$compute(i_1)$};
\node[inner, below=of c1] (u1) {$update(i_1)$};

\node[bott, fit=(r1)(c1)(u1), label=above:{$F(i_1)$}] (F1) {};

\draw[arrow] (r1) -- (c1);
\draw[arrow] (c1) -- (u1);

\node[below=10mm of F1] (dots) {$\vdots$};

\node[inner, below=10mm of dots] (rn) {$read(i_n)$};
\node[inner, below=of rn] (cn) {$compute(i_n)$};
\node[inner, below=of cn] (un) {$update(i_n)$};

\node[bott, fit=(rn)(cn)(un), label=above:{$F(i_n)$}] (Fn) {};

\draw[arrow] (rn) -- (cn);
\draw[arrow] (cn) -- (un);

\draw[arrow] (F1.south) -- node[right] {secuencial} (dots);
\draw[arrow] (dots) -- (Fn.north);

\node[above=6mm of F1] {\textbf{Estado \(P\):} \(P = F(i_1);F(i_2);\dots;F(i_n)\)};

\end{tikzpicture}

\caption{Pipeline secuencial del sistema.}
\end{figure}

\section{Optimización del sistema}

A partir del análisis del Estado \(P\), se observa que el principal problema del sistema radica en la ejecución estrictamente secuencial del pipeline de procesamiento. Esta estructura impide el aprovechamiento de recursos computacionales disponibles, particularmente en sistemas modernos con múltiples núcleos de procesamiento.

Con el objetivo de mejorar el rendimiento del sistema, se analiza la naturaleza de las tareas que componen el pipeline para identificar cuáles pueden beneficiarse de estrategias de concurrencia.

\subsection{Identificación de tareas I/O-bound}

Las tareas \textit{I/O-bound} son aquellas cuyo tiempo de ejecución está dominado por operaciones de entrada/salida, tales como acceso a sensores, lectura de datos o comunicación con dispositivos externos.

En el sistema propuesto, la operación:

\[
read(i)
\]

representa la adquisición de datos de tráfico provenientes de sensores o cámaras ubicadas en cada intersección.

Formalmente:

\[
read : I \rightarrow D
\]

donde \(D\) representa el conjunto de datos de tráfico capturados por los sensores.

Estas operaciones suelen presentar latencias asociadas al hardware o a la comunicación con dispositivos externos, por lo que múltiples operaciones de lectura pueden ejecutarse de forma concurrente sin interferencia entre sí.

\subsection{Identificación de tareas CPU-bound}

Las tareas \textit{CPU-bound} son aquellas cuyo tiempo de ejecución está dominado por cálculos computacionales.

En el pipeline definido, la operación:

\[
compute(i)
\]

corresponde al cálculo del tiempo óptimo del semáforo en la intersección \(i\), considerando variables como:

\begin{itemize}
\item densidad vehicular
\item flujo de tráfico
\item tiempos de espera
\end{itemize}

Formalmente:

\[
compute : D \rightarrow T
\]

donde \(T\) representa el tiempo calculado para el ciclo del semáforo.

Estas operaciones requieren procesamiento intensivo y pueden beneficiarse de paralelización mediante múltiples procesos o hilos de ejecución.

\subsection{Estrategia de paralelización}

Dado que las intersecciones de tráfico pueden procesarse de manera independiente, es posible distribuir el procesamiento de múltiples pipelines entre diferentes unidades de ejecución.

Sea:

\[
F(i) = read(i); compute(i); update(i)
\]

En el modelo secuencial:

\[
P = F(i_1); F(i_2); \dots ; F(i_n)
\]

En contraste, una estrategia de paralelización permite ejecutar múltiples pipelines simultáneamente:

\[
Q = \parallel_{i=1}^{n} F(i)
\]

donde el operador \(\parallel\) representa ejecución concurrente.

Esta transformación permite reducir el tiempo total de ejecución del sistema.

\subsection{Estado Q: Sistema Concurrente}

El Estado \(Q\) representa una versión optimizada del sistema en la cual múltiples intersecciones pueden procesarse simultáneamente.

Formalmente:

\[
Q = \parallel_{i=1}^{n} (read(i); compute(i); update(i))
\]

El tiempo de ejecución aproximado se convierte en:

\[
T_Q \approx \max (t_r + t_c + t_u)
\]

asumiendo una distribución equilibrada de las tareas entre los recursos disponibles.

Desde la perspectiva de complejidad paralela:

\[
T_Q = O\left(\frac{n}{p}\right)
\]

donde \(p\) representa el número de unidades de procesamiento disponibles.

La transformación del sistema puede expresarse formalmente mediante una tripleta de Hoare:

\[
\{P\} \ OptimizeConcurrency \ \{Q\}
\]

lo cual indica que el proceso de optimización transforma el sistema secuencial en una versión concurrente manteniendo la corrección funcional.

\subsection{Complejidad del sistema concurrente}

El modelo concurrente del sistema, representado por el Estado \(Q\), permite ejecutar múltiples pipelines de procesamiento de forma simultánea. A diferencia del Estado \(P\), donde cada intersección debe esperar a que la anterior finalice su ejecución, en el modelo concurrente las intersecciones pueden procesarse de manera independiente.

Recordando que el pipeline por intersección se define como:

\[
F(i) = read(i); compute(i); update(i)
\]

y que el costo computacional de cada pipeline es constante:

\[
F(i) = O(1)
\]

el sistema secuencial se modelaba como:

\[
P = F(i_1); F(i_2); \dots ; F(i_n)
\]

lo cual produce una complejidad:

\[
T_P = O(n)
\]

En el modelo concurrente, el sistema se transforma en:

\[
Q = \parallel_{i=1}^{n} F(i)
\]

donde el operador \(\parallel\) representa ejecución paralela de múltiples pipelines.

Si el sistema dispone de \(p\) unidades de procesamiento, el tiempo de ejecución aproximado puede expresarse como:

\[
T_Q = O\left(\frac{n}{p}\right)
\]

donde:

\begin{itemize}
\item \(n\) es el número de intersecciones
\item \(p\) es el número de unidades de procesamiento disponibles
\end{itemize}

Este modelo indica que el tiempo de ejecución del sistema puede reducirse proporcionalmente al número de recursos computacionales utilizados, siempre que el trabajo pueda distribuirse de manera equilibrada.

\subsection{Speedup esperado}

El rendimiento del sistema concurrente puede evaluarse mediante el concepto de \textit{speedup}, definido como:

\[
S = \frac{T_P}{T_Q}
\]

Sustituyendo los modelos obtenidos:

\[
S = \frac{O(n)}{O(n/p)}
\]

lo cual produce:

\[
S \approx p
\]

Este resultado indica que, en condiciones ideales, el sistema concurrente puede alcanzar un speedup cercano al número de procesadores disponibles.

Sin embargo, en la práctica existen factores que limitan el rendimiento, tales como:

\begin{itemize}
\item sincronización entre procesos
\item sobrecarga de creación de hilos o procesos
\item desbalance de carga
\end{itemize}

Por lo tanto, el speedup real del sistema dependerá de la eficiencia de la estrategia de paralelización utilizada.

\subsection{Diagrama del pipeline concurrente}

El Estado \(Q\) permite que múltiples pipelines de procesamiento se ejecuten simultáneamente, eliminando la dependencia estricta entre intersecciones.

\begin{figure}[H]
\centering
\begin{tikzpicture}[
block/.style={rectangle, draw, rounded corners, align=center, minimum width=2.5cm, minimum height=7mm},
arrow/.style={-Latex, thick}
]

\node[block] (r1) {read($i_1$)};
\node[block, below=of r1] (c1) {compute($i_1$)};
\node[block, below=of c1] (u1) {update($i_1$)};

\node[block, right=3.5cm of r1] (r2) {read($i_2$)};
\node[block, below=of r2] (c2) {compute($i_2$)};
\node[block, below=of c2] (u2) {update($i_2$)};

\node[right=3cm of r2] (dots) {$\cdots$};

\node[block, right=3cm of dots] (rn) {read($i_n$)};
\node[block, below=of rn] (cn) {compute($i_n$)};
\node[block, below=of cn] (un) {update($i_n$)};

\draw[arrow] (r1) -- (c1);
\draw[arrow] (c1) -- (u1);

\draw[arrow] (r2) -- (c2);
\draw[arrow] (c2) -- (u2);

\draw[arrow] (rn) -- (cn);
\draw[arrow] (cn) -- (un);

\end{tikzpicture}

\caption{Pipeline concurrente del sistema. Cada intersección se procesa de forma independiente.}
\end{figure}

\end{document}